\cotefile{135}
\scanpage{48}
\mspage{2}
of bifunctors
\[
c_{X,Y}:X\otimes Y\xrightarrow{\sim}Y\otimes X,\qquad X,Y\in\operatorname{Ob}\underline{C}.
\]
verifying the relation
\[
c_{Y,X}\circ c_{X,Y}=\mathrm{id}_{X\otimes Y}
\]
The commutativity constraint $c$ is said to be \uline{strict} if $c_{X,X}=\mathrm{id}_{X\otimes X}$ for all $X\in\operatorname{Ob}\underline{C}$. A $\otimes$-category together with a commutativity constraint is a \uline{$\otimes$-commutative category}. A $\otimes$-commutative category is \uline{strict} if its commutativity constraint is strict.

An \uline{unity constraint} for a $\otimes$-category $\underline{C}$ is a triple $(\underline{1},g,d)$ where $\underline{1}$ is an object of $\underline{C}$, $g$ and $d$ natural isomorphisms
\[
g_X:X\xrightarrow{\sim}\underline{1}\otimes X,\qquad
d_X:X\xrightarrow{\sim}X\otimes\underline{1},\qquad X\in\operatorname{Ob}\underline{C}
\]
such that $g_{\underline{1}}=d_{\underline{1}}$. A $\otimes$-category together with an unity constraint is a \uline{$\otimes$-unifer category}.

A $\otimes$-category $\underline{C}$ together with an associativity constraint $a$ and a commutativity constraint $c$ is a \uline{$\otimes$-AC category} if the \uline{hexagon axiom} is fulfilled, i.e all the hexagonal diagrams commute
\diagram{\begin{tikzcd}[ampersand replacement=\&,column sep=small,row sep=large]
\& (X\otimes Y)\otimes Z \arrow[r,"{c_{X\otimes Y,Z}}"] \& Z\otimes(X\otimes Y) \arrow[dr,"{a_{Z,X,Y}}"] \& \\
X\otimes(Y\otimes Z) \arrow[ur,"{a_{X,Y,Z}}"] \arrow[dr,"{\mathrm{id}_X\otimes c_{Y,Z}}"'] \& \& \& (Z\otimes X)\otimes Y\\
\& X\otimes(Z\otimes Y) \arrow[r,"{a_{X,Z,Y}}"] \& (X\otimes Z)\otimes Y \arrow[ur,"{c_{X,Z}\otimes\mathrm{id}_Y}"'] \&
\end{tikzcd}}

A $\otimes$-category $\underline{C}$ together with an associativity constraint $a$ and an unity constraint $(\underline{1},g,d)$ is a \uline{$\otimes$-AU category} if all the following triangles commute
\diagram{\begin{tikzcd}[ampersand replacement=\&,column sep=small]
X\otimes(\underline{1}\otimes Y) \arrow[rr,"{a_{X,\underline{1},Y}}"] \& \& (X\otimes\underline{1})\otimes Y\\
\& X\otimes Y \arrow[ul,"{\mathrm{id}_X\otimes g_Y}"] \arrow[ur,"{d_X\otimes\mathrm{id}_Y}"'] \&
\end{tikzcd}}
% archive sheet number: 47
\scanpage{49}
\mspage{3}
A \uline{$\otimes$-ACU category} is a $\otimes$-AC and AU category. An object $X$ of a $\otimes$-ACU category $\underline{C}$ is \uline{invertible} if there are two objects $X',X''\in\operatorname{Ob}\underline{C}$ such that $X'\otimes X\simeq X\otimes X''\simeq\underline{1}$.

A \uline{$\otimes$-functor} from a $\otimes$-category $\underline{C}$ to a $\otimes$-category $\underline{C}'$ is a pair $(F,\breve F)$ where $F$ is a functor $\underline{C}\longrightarrow\underline{C}'$ and $\breve F$ an isomorphism of bifunctors
\[
\breve F_{X,Y}:FX\otimes FY\longrightarrow F(X\otimes Y)\qquad X,Y\in\operatorname{Ob}\underline{C}
\]

A $\otimes$-functor $(F,\breve F)$ from a $\otimes$-associative category $\underline{C}$ to a $\otimes$-associative \unsure{category|catégory} $\underline{C}'$ is \uline{associative} if the following diagram commutes:
\diagram{\begin{tikzcd}[ampersand replacement=\&,column sep=small]
FX\otimes(FY\otimes FZ) \arrow[r,"{\mathrm{id}\otimes\breve F}"] \arrow[d,"{a'}"'] \& FX\otimes F(Y\otimes Z) \arrow[r,"{\breve F}"] \& F(X\otimes(Y\otimes Z)) \arrow[d,"{Fa}"]\\
(FX\otimes FY)\otimes FZ \arrow[r,"{\breve F\otimes\mathrm{id}}"] \& F(X\otimes Y)\otimes FZ \arrow[r,"{\breve F}"] \& F((X\otimes Y)\otimes Z)
\end{tikzcd}}
where $a$ is the associativity constraint of $\underline{C}$ and $a'$ of $\underline{C}'$.

A $\otimes$-functor $(F,\breve F)$ from a $\otimes$-commutative category $\underline{C}$ to a $\otimes$-commutative category $\underline{C}'$ is \uline{commutative} if the following diagram commutes
\diagram{\begin{tikzcd}[ampersand replacement=\&]
FX\otimes FY \arrow[r,"{\breve F}"] \arrow[d,"{c'}"'] \& F(X\otimes Y) \arrow[d,"{Fc}"]\\
FY\otimes FX \arrow[r,"{\breve F}"] \& F(Y\otimes X)
\end{tikzcd}}
$c$ and $c'$ being the commutativity constraints of $\underline{C}$ and $\underline{C}'$ respectively.

A $\otimes$-functor $(F,\breve F)$ from a $\otimes$-category $\underline{C}$ with an unity constraint $(\underline{1},g,d)$ to a $\otimes$-category $\underline{C}'$ with an unity constraint $(\underline{1}',g',d')$ is a \uline{$\otimes$-unifer functor} if there exists an isomorphism $\hat F:\underline{1}'\xrightarrow{\sim}F\underline{1}$ such that the following diagrams commute:
\diagram{\begin{tikzcd}[ampersand replacement=\&,column sep=large]
\underline{1}'\otimes FX \arrow[r,"{\hat F\otimes\mathrm{id}_{FX}}"] \& F\underline{1}\otimes FX \arrow[d,"{\breve F}"]\\
FX \arrow[u,"{g'_{FX}}"] \arrow[r,"{Fg_X}"'] \& F(\underline{1}\otimes X)
\end{tikzcd}\qquad
\begin{tikzcd}[ampersand replacement=\&,column sep=large]
FX\otimes\underline{1}' \arrow[r,"{\mathrm{id}_{FX}\otimes\hat F}"] \& FX\otimes F\underline{1} \arrow[d,"{\breve F}"]\\
FX \arrow[u,"{d'_{FX}}"] \arrow[r,"{Fd_X}"'] \& F(X\otimes\underline{1})
\end{tikzcd}}
% archive sheet number: 48
\scanpage{50}
\mspage{4}
It follows from the definition that the isomorphism $\hat F:\underline{1}'\xrightarrow{\sim}F\underline{1}$, if it exists, is unique.

A \uline{$\otimes$-AC functor} is \unsure{an|a} $\otimes$-associative and commutative functor. A \uline{$\otimes$-ACU functor} is a $\otimes$-associative, commutative and unifer functor.

Let $(F,\breve F)$ and $(G,\breve G)$ be $\otimes$-functors from a $\otimes$-category $\underline{C}$ to a $\otimes$-category $\underline{C}'$. A \uline{$\otimes$-morphism} from the $\otimes$-functor $(F,\breve F)$ to the $\otimes$-functor $(G,\breve G)$ is a morphism of functors $\lambda:F\to G$ such that the following diagram commutes
\diagram{\begin{tikzcd}[ampersand replacement=\&]
FX\otimes FY \arrow[r,"{\breve F}"] \arrow[d,"{\lambda_X\otimes\lambda_Y}"'] \& F(X\otimes Y) \arrow[d,"{\lambda_{X\otimes Y}}"]\\
GX\otimes GY \arrow[r,"{\breve G}"] \& G(X\otimes Y)
\end{tikzcd}\qquad $X,Y\in\operatorname{Ob}\underline{C}$}

Chapter $\overline{\underline{\mathrm{II}}}$ is a study of Gr-categories and Pic-categories. A \uline{Gr-category} is a $\otimes$-AU category, the objects of which are all invertible, and the base category a groupoïd (i.e all arrows are isomorphisms). Thus a Gr-category is like a group. We obtain from this definition that if $\underline{P}$ is a Gr-category, the set $\Pi_0(\underline{P})$ of the classes up to ismomorphism of objects of $\underline{P}$, together with the operation induced by the law $\otimes$ of $\underline{P}$, is a group; the group $\operatorname{Aut}(\underline{1})=\Pi_1(\underline{P})$ is a commutative group; and for all $X\in\operatorname{Ob}\underline{P}$
\[
\begin{aligned}
\gamma_X:\ u\longmapsto u\otimes\mathrm{id}_X&=\operatorname{Aut}(\underline{1})\xrightarrow{\sim}\operatorname{Aut}(X)\\
\delta_X:\ u\longmapsto\mathrm{id}_X\otimes u&=\operatorname{Aut}(\underline{1})\xrightarrow{\sim}\operatorname{Aut}(X)
\end{aligned}
\]
We attribute thus to a Gr-category $\underline{P}$ two groups $\Pi_0(\underline{P})$ and $\Pi_1(\underline{P})$ where $\Pi_1(\underline{P})$ is commutative. Furthermore we can define an action of $\Pi_0(\underline{P})$ on $\Pi_1(\underline{P})$ by the formula
\[
su=\delta_X^{-1}\gamma_X(u)
\]
for $s\in\Pi_0(\underline{P})$ represented by $X$ and $u\in\Pi_1(\underline{P})$. The commutative group $\Pi_1(\underline{P})$ together with this action is a left $\Pi_0(\underline{P})$-module.

Let $M$ be a group, $N$ a left $M$-module. A \uline{preepinglages} of
% archive sheet number: 49
\scanpage{51}
\mspage{5}
of type $(M,N)$ for a Gr-category $\underline{P}$ is a pair $\varepsilon=(\varepsilon_0,\varepsilon_1)$ of isomorphisms
\[
\varepsilon_0:M\xrightarrow{\sim}\Pi_0(\underline P),\qquad
\varepsilon_1:N\xrightarrow{\sim}\Pi_1(\underline P).
\]
compatible with the action of $M$ on $N$, $\Pi_0(\underline P)$ on $\Pi_1(\underline P)$. A Gr-category \margin{!}\uline{pree\mbox{\tikz[baseline=(p.base)]{\node[inner sep=0pt] (p) {ping};\draw[rotate around={28:(p.center)}] (p.center) ellipse [x radius=.9em,y radius=.72em];}}led} of type $(M,N)$ is a Gr-category $\underline P$ together with a preepinglage. Finally, an \uline{arrow} of Gr-categories preepingled of type $(M,N)$ $(\underline P,\varepsilon)\to(\underline P',\varepsilon')$ is a $\otimes$-associative functor such that the following triangles commute:
\diagram{\begin{tikzcd}[ampersand replacement=\&,column sep=small]
\Pi_0(\underline P) \arrow[rr] \& \& \Pi_0(\underline P')\\
\& M \arrow[ul,"{\varepsilon_0}"] \arrow[ur,"{\varepsilon'_0}"'] \&
\end{tikzcd}\qquad
\begin{tikzcd}[ampersand replacement=\&,column sep=small]
\Pi_1(\underline P) \arrow[rr] \& \& \Pi_1(\underline P')\\
\& N \arrow[ul,"{\varepsilon_1}"] \arrow[ur,"{\varepsilon'_1}"'] \&
\end{tikzcd}}

It follows from this definition that a such arrow is a $\otimes$-equivalence. Thus the set of the equivalence classes of Gr-categories preepingled of type $(M,N)$ is equal to the set of the connected components of the category of Gr-categories preepingled of type $(M,N)$.

If we consider the cohomology group $H^3(M,N)$ of the group $M$ with coefficients $N$ (in the sense of the group cohomoly \unsure{.|{}} [12]) we obtain a canonical bijection between the set $H^3(M,N)$ and the set of the equivalence classes of Gr-categories preepingled of type $(M,N)$.

A \uline{Pic-category} is a Gr-category together with a commutativity constraint which is compatible with its associativity constraint, i.e the hexagon axiom is satisfied. Thus a Pic-category is like a commutative group. We verify immediately that a necessary condition for the existence of a Pic-category structure on a Gr-category is that $\Pi_0(\underline P)$ must be commutative and act trivially on $\Pi_1(\underline P)$. A Pic-category is \uline{strict} if its commutativity constraint is strict.

\unsure{let|Let} $M,N$ be abelian groups. A \uline{preepinglage} of type $(M,N)$ for a Pic-category $\underline P$ is a pair $\varepsilon=(\varepsilon_0,\varepsilon_1)$ of isomorphisms
\[
\varepsilon_0:M\xrightarrow{\sim}\Pi_0(\underline P),\qquad
\varepsilon_1:N\xrightarrow{\sim}\Pi_1(\underline P)
\]
% archive sheet number: 50
\scanpage{52}
\mspage{6}
A Pic-category \uline{preepingled} of type $(M,N)$ is a Pic-category together with a preepinglage. We define the \uline{arrow} of such objects in the same way as for Gr-categories.

For next propositions, let us consider two complexes of free abelian groups
\[
\begin{aligned}
L_{\unsure{\cdot|\bullet|.}}(M):\quad &L_3(M)\xrightarrow{d_3}L_2(M)\xrightarrow{d_2}L_1(M)\xrightarrow{d_1}L_0(M)\longrightarrow M\\
{}'L_{\cdot}(M):\quad &{}'L_3(M)\xrightarrow{{}'d_3}{}'L_2(M)\xrightarrow{{}'d_2}{}'L_1(M)\xrightarrow{{}'d_1}{}'L_0(M)\longrightarrow M
\end{aligned}
\]
where
\[
\resizebox{\linewidth}{!}{$\displaystyle
\begin{aligned}
L_0(M)&={}'L_0(M)=\mathbb Z[M]\\
L_1(M)&={}'L_1(M)=\mathbb Z[M\times M]\\
L_2(M)&={}'L_2(M)=\mathbb Z[M\times M\times M]+\mathbb Z[M\times M]\\
L_3(M)&={}'L_3(M)+\mathbb Z[M]\\
{}'L_3(M)&=\mathbb Z[M\times M\times M\times M]+\mathbb Z[M\times M\times M]+\mathbb Z[M\times M]\\
d_1[x,y]&={}'d_1[x,y]=[y]-[x+y]+[x]\\
d_2[x,y]&={}'d_2[x,y]=[x,y]-[y,x]\\
d_2[x,y,z]&={}'d_2[x,y,z]=[y,z]-[x+y,z]+[x,y+z]-[x,y]\\
d_3[x,y,z,t]&={}'d_3[x,y,z,t]=[y,z,t]-[x+y,z,t]+[x,y+z,t]\\
&\hspace{11em}{}-[x,y,z+t]+[x,y,z]\\
d_3[x,y,z]&={}'d_3[x,y,z]=[x,y,z]-[x,z,y]+[z,x,y]-[y\unsure{\mathord{\cdot}|{,}}z]+{}\\
&\hspace{9em}{}+[x+y,z]-[x,z].\\
d_3[x,y]&=[x,y]+[y,x]={}'d_3[x,y]\\
d_3[x]&=[x,x],
\end{aligned}$}
\]
so that $L_{\cdot}(M)$ is a truncated resolution of $M$. One obtains a canonical bijection between the set of the equivalence classes of
% archive sheet number: 51
\scanpage{53}
\mspage{7}
Pic-categories preepingled of type $(M,N)$ and the set $H^2(\operatorname{Hom}({}'L_{\unsure{\cdot|\bullet|.}}(M),N))$
The exactitude of the complex $L_{\cdot}(M)$ gives us the triviality of the classification of Pic-categories preepingled of type $(M,N)$ which are strict, i.e all Pic-categories preepingled of type $(M,N)$ which are strict, are equivalent.
\margin{\unsure{Determine|Déterminer} $H^2(\operatorname{Hom}({}'L_{\cdot}(M),N))$ i.e. the}

Finally chapter $\overline{\underline{\mathrm{III}}}$ gives us the construction of the solution of two universal problems: \uline{problem of making objects ``unity \unsure{objects|object}''} and \uline{problem of reversing objects}.

Let $\underline A$ be a $\otimes$-AC category, $\underline A'$ another $\otimes$-AC category whose base category is a groupoïd, and $(T,\breve T):\underline A'\to\underline A$ a $\otimes$-AC functor. We try to make the objects $TA'$ of $\underline A$, $A'\in\operatorname{Ob}\underline A'$, ``unity object'',
% slanted
\margin{\unsure{simply|strictly|auxiliary} unit?}
i.e we try to get:

$1^\circ$ A $\otimes$-ACU category $\underline P$

$2^\circ$ A $\otimes$-AC functor $(D,\breve D):\underline A\to\underline P$

$3^\circ$ A $\otimes$-isomorphism
\[
\lambda:(D,\breve D)\circ(T,\breve T)\xrightarrow{\sim}(I_{\underline P},\breve I_{\underline P}).
\]
where $(I_{\underline P},\breve I_{\underline P})$ is the $\otimes$-constant functor $\underline{1}_{\underline P}$ from $\underline A'$ to $\underline P$. The triple $(\underline P,(D,\breve D),\lambda)$ must be universal for triples $(\underline Q,(E,\breve E),\mu)$ satisfying $1^\circ,2^\circ,3^\circ$.

For the description of the triple $(\underline P,(D,\breve D),\lambda)$, we introduce \margin{$\big|$}a quotient category of a $\otimes$-AC category as follows:

Let $\underline A$ be a $\otimes$-AC category, $\unsure{\mathscr S|\mathcal S}$ a \uline{multiplicative subset} of $\underline A$ (that means a subset of the set of all \deleted{endomorphisms}\inserted{arrows} of $\underline A$
% slanted
\margin{what about composition?}
such that $\mathrm{id}_X\in\mathscr S$ for all $X\in\operatorname{Ob}\underline A$ and the tensor product of two arrows of $\mathscr S$ belongs to $\mathscr S$). The \uline{$\otimes$-AC \mbox{\tikz[baseline=(w.base)]{\node[inner sep=0pt] (w) {category quotient};\draw (-.25,.19) .. controls (.13,.20) and (.04,.08) .. (.06,0) .. controls (.06,-.10) and (.15,-.11) .. (.43,-.11);}}} $A^{\mathscr S}$ of $\underline A$ with respect to $\mathscr S$ is the solution of the universal problem
\[
(K,\breve K):\underline A\longrightarrow\underline B,\qquad K(u)=\mathrm{id}\quad\text{for all }u\in\mathscr S
\]
where $B$ is a $\otimes$-AC category and $(K,\breve K)$ a $\otimes$-AC functor.
% archive sheet number: 52
\scanpage{54}
\mspage{8}
Now let us give an idea of the construction of the triple $(\underline P,(D,\breve D),\lambda)$ for $\underline A'\ne\varnothing$:

\[
\begin{aligned}
1^\circ\quad &\operatorname{Ob}\underline P=\operatorname{Ob}\underline A\\
2^\circ\quad &\operatorname{Hom}_{\underline P}(A,B)=\unsure{\phi|\Phi}(A,B)/_{\unsure{\mathcal R|\mathscr R}_{A,B}},\qquad A,B\in\operatorname{Ob}\underline P
\end{aligned}
\]
$\phi(A,B)$ being the set of all triples $(A',B',u)$ where $A',B'\in\operatorname{Ob}\underline A'$, $u\in\operatorname{Fl}\underline A$ $u:A\otimes TA'\longrightarrow B\otimes TB'$; $\mathcal R_{A,B}$ the equivalence relation defined in $\phi(A,B)$ as follows
\[
(A'_1,B'_1,u)\ \mathcal R_{A,B}\ (A'_2,B'_2,u)
\]
if and only if there are objects $C'_1,C'_2$ and isomorphisms
\[
u':A'_1\otimes C'_1\xrightarrow{\sim}A'_2\otimes C'_2,\qquad v':B'_1\otimes C'_1\xrightarrow{\sim}B'_2\otimes C'_2
\]
of $\underline A'$ such that the following diagram commutes in $\underline A^{\unsure{\mathscr S|\mathcal S}}$, $\otimes$-AC quotient category of $\underline A$ with respect to the multiplicative subset of $\underline A$ generated by the endomorphisms of the form $T(c_{A',A'})$;
\diagram{\resizebox{\linewidth}{!}{\begin{tikzcd}[ampersand replacement=\&,column sep=large,row sep=large]
A\otimes T(A'_1\otimes C'_1) \arrow[r,"{\mathrm{id}\otimes\breve T^{-1}}"] \arrow[d,"{\mathrm{id}\otimes Tu'}"] \& A\otimes(TA'_1\otimes TC'_1) \arrow[r,"{a}"] \& (A\otimes TA'_1)\otimes TC'_1 \arrow[r,"{u_1\otimes\mathrm{id}}"] \& (B\otimes TB'_1)\otimes{TC'_1} \arrow[d,"{a^{-1}}"']\\
A\otimes T(A'_2\otimes C'_2) \arrow[d,"{\mathrm{id}\otimes\breve T^{-1}}"] \& \& \& B\otimes(TB'_1\otimes TC'_1) \arrow[d,"{\mathrm{id}\otimes\breve T}"']\\
A\otimes(TA'_2\otimes TC'_2) \arrow[d,"{a}"] \& \& \& B\otimes T(B'_1\otimes C'_1) \arrow[d,"{\mathrm{id}\otimes Tv'}"']\\
(A\otimes TA'_2)\otimes TC'_2 \arrow[r,"{u_2\otimes\mathrm{id}}"] \& (B\otimes TB'_2)\otimes TC'_2 \arrow[r,"{a^{-1}}"] \& B\otimes(TB'_2\otimes TC'_2) \arrow[r,"{\mathrm{id}\otimes\breve T^{-1}}"] \& B\otimes T(B'_2\otimes C'_2)
\end{tikzcd}}}
We denote by $[A',B',u]$ the class which has $(A',B',u)$ as representative

$3^\circ$ Composition of arrows in $\underline P$. Let $[A',B',u]:A\to B$, $[B'',C'',v]:B\to C$ be arrows in $\underline P$. We define
\[
[B'',C'',v]\circ[A',B',u]=[A'\otimes B'',B'\otimes C'',\omega]:A\longrightarrow C
\]
% archive sheet number: 53
\scanpage{55}
\mspage{9}
where $\omega$ is such that the following diagram commutes:
\diagram{\resizebox{\linewidth}{!}{\begin{tikzcd}[ampersand replacement=\&,column sep=large,row sep=small]
A\otimes T(A'\otimes B'') \arrow[r,"{\mathrm{id}\otimes\breve T^{-1}}"] \arrow[dddd,"{\omega}"] \& A\otimes(TA'\otimes TB'') \arrow[r,"{a}"] \& (A\otimes TA')\otimes TB'' \arrow[r,"{u\otimes\mathrm{id}}"] \& (B\otimes TB')\otimes TB'' \arrow[d,"{a^{-1}}"]\\
\& \& \& B\otimes(TB'\otimes TB'') \arrow[d,"{\mathrm{id}\otimes c}"]\\
\& \& \& B\otimes(TB''\otimes TB') \arrow[d,"{a}"]\\
\& \& \& (B\otimes TB'')\otimes TB' \arrow[d,"{v\otimes\mathrm{id}}"]\\
C\otimes T(B'\otimes C'') \& C\otimes(TB'\otimes TC'') \arrow[l,"{\mathrm{id}\otimes\breve T}"'] \& C\otimes(TC''\otimes TB') \arrow[l,"{\mathrm{id}\otimes c}"'] \& (C\otimes TC'')\otimes TB' \arrow[l,"{a^{-1}}"']
\end{tikzcd}}}

$4^\circ$ $\otimes$-Structure on $\underline P$
\[
\begin{aligned}
A\otimes E\ (\text{in }\underline P)&=A\otimes E\ (\text{in }\underline A)\\
[A',B',u]\otimes[E',F',v]&=[A'\otimes E',B'\otimes F',w]
\end{aligned}
\]
where $w$ is defined by the commutative diagram (1).

$5^\circ$ ACU constraint in $\underline P$
\[
\resizebox{\linewidth}{!}{$\displaystyle
\bigl([A',A',a\otimes\mathrm{id}],\ [A',A',c\otimes\mathrm{id}],\
(\underline{1}_{\underline P}=TA'_0,\ g_A=[A'_0\otimes A',A',t_A],\
d_A=[A'_0\otimes A',A',\unsure{p|\rho|r}_A])\bigr)$}
\]
where $A'_0$ is a fixed object\unsure{,|{}} of $\underline A'$, $A'$ an arbitrary object of $\underline A'$, $g_A$ and $d_A$ natural isomorphisms
\[
g_A:A\longrightarrow\underline{1}_{\underline P}\otimes A,\qquad
d_A:A\longrightarrow A\otimes\underline{1}_{\underline P}
\]
with $t_A$ and $p_A$ defined by the commutative diagrams (2)

$6^\circ$ $(D,\breve D)$ is defined by
\[
DA=A,\qquad Du=[A',A',u\otimes\mathrm{id}_{TA'}],\qquad\breve D_{A,B}=\mathrm{id}_{A\otimes B}
\]
% archive sheet number: 54
\scanpage{56}
\mspage{10}
\diagram{\resizebox{\linewidth}{!}{\begin{tikzcd}[ampersand replacement=\&,column sep=large,row sep=large]
\& (A\otimes TA')\otimes(E\otimes TE') \arrow[r,"{u\otimes v}"] \arrow[d,"{a}"'] \& (B\otimes TB')\otimes(F\otimes TF') \arrow[d,"{a}"]\\
\& ((A\otimes TA')\otimes E)\otimes TE' \arrow[d,"{a^{-1}\otimes\mathrm{id}}"'] \& ((B\otimes TB')\otimes F)\otimes TF' \arrow[d,"{a^{-1}\otimes\mathrm{id}}"]\\
\& (A\otimes(TA'\otimes E))\otimes TE' \arrow[d,"{(\mathrm{id}\otimes c)\otimes\mathrm{id}}"'] \& (B\otimes(TB'\otimes F))\otimes TF' \arrow[d,"{(\mathrm{id}\otimes c)\otimes\mathrm{id}}"]\\
(1) \& (A\otimes(E\otimes TA'))\otimes TE' \arrow[d,"{a\otimes\mathrm{id}}"'] \& (B\otimes(F\otimes TB'))\otimes TF' \arrow[d,"{a\otimes\mathrm{id}}"]\\
\& ((A\otimes E)\otimes TA')\otimes TE' \arrow[d,"{a^{-1}}"'] \& ((B\otimes F)\otimes TB')\otimes TF' \arrow[d,"{a^{-1}}"]\\
\& (A\otimes E)\otimes(TA'\otimes TE') \arrow[d,"{\mathrm{id}\otimes\breve T}"'] \& (B\otimes F)\otimes(TB'\otimes TF') \arrow[d,"{\mathrm{id}\otimes\breve T}"]\\
\& (A\otimes E)\otimes T(A'\otimes E') \arrow[r,"{w}"] \& (B\otimes F)\otimes T(B'\otimes F')
\end{tikzcd}}}

\diagram{\resizebox{\linewidth}{!}{(2)\quad
\begin{tikzcd}[ampersand replacement=\&,column sep=large]
A\otimes(TA'_0\otimes TA') \arrow[r,"{\mathrm{id}\otimes\breve T}"] \arrow[d,"{a}"'] \& A\otimes T(A'_0\otimes A') \arrow[d,"{t_A}"]\\
(A\otimes TA'_0)\otimes TA' \arrow[r,"{c\otimes\mathrm{id}}"] \& (TA'_0\otimes A)\otimes TA'
\end{tikzcd}\qquad
\begin{tikzcd}[ampersand replacement=\&,column sep=large]
A\otimes(TA'_0\otimes TA') \arrow[r,"{\mathrm{id}\otimes\breve T}"] \arrow[d,"{a}"'] \& A\otimes T(A'_0\otimes A') \arrow[d,"{\unsure{p|\rho|r}_A}"]\\
(A\otimes TA'_0)\otimes TA' \arrow[r,equal] \& (A\otimes TA'_0)\otimes TA'
\end{tikzcd}}}

$7^\circ$ The $\otimes$-isomorphism
\[
\lambda:(D,\breve D)\circ(T,\breve T)\xrightarrow{\sim}(I_{\underline P},\breve I_{\underline P})
\]
is defined by natural isomorphisms
\[
DTA'=TA'\xrightarrow{\lambda_{A'}=[A'_0,A',c_{TA',TA'_0}]}I_{\underline P}A'=TA'_0\qquad A'\in\operatorname{Ob}\underline A'
\]

\margin{?}$\underline P$ is called the \uline{$\otimes$-A\mbox{\tikz[baseline=(w.base)]{\node[inner sep=0pt] (w) {CU category};\draw[rotate around={6:(w.center)}] (w.center) ellipse [x radius=2.8em,y radius=.73em];}}} of the $\otimes$-AC category $\underline A$ with respect to $(A',(T,\breve T))$. \margin{examples?}

For the problem of reversing objects, let us consider a $\otimes$-category
% archive sheet number: 55
\scanpage{57}
\mspage{11}
$\underline C$ with a ACU constraint $(a,c,(\underline{1},g,d))$, a $\otimes$-category $\underline C'$ with a ACU constraint $(a',c',(\underline{1}',g',d'))$, the base category of which is a groupoïd, and a $\otimes$-ACU functor $(F,\breve F):\underline C'\longrightarrow\underline C$. We try to find a $\otimes$-ACU category $\underline P$ and a $\otimes$-ACU functor $(\unsure{\mathfrak D|\mathscr D|Q},\breve{\mathfrak D}):\underline C\longrightarrow\underline P$ having the following properties:

$1^\circ$ $\mathfrak DFX'$ is invertible in $\underline P$ for all $X'\in\operatorname{Ob}\underline C'$.

$2^\circ$ For all $\otimes$-ACU functor $(\unsure{\mathfrak E|\mathscr E|Y},\breve{\mathfrak E})$ from $\underline C$ to a $\otimes$-ACU category $\underline Q$
% slanted
\margin{\unsure{give} the \unsure{equivalence|equivalency|équivalence} of categories \tikz[baseline]{\draw (0,.55) -- (.025,-.28) .. controls (.025,-.31) and (-.015,-.31) .. (-.025,-.28);}}
such that $\mathfrak EFX'$ is invertible in $\underline Q$ for all $X'\in\operatorname{Ob}\underline C'$, there exists a $\otimes$-ACU functor $(E',\breve E')$, unique up to $\otimes$-isomorphism, from $\underline P$ to $\underline Q$ such that $(\mathfrak E,\breve{\mathfrak E})\simeq(E',\breve E')\circ(\mathfrak D,\breve{\mathfrak D})$.

This problem is reduced by the first by putting $\underline A'=\underline C'$, $\underline A=\underline C\times\underline C'$, $TX'=(FX',X')$ and by remarking that if $\underline C,\underline C',\underline Q$ are $\otimes$-ACU categories, $\underline{\operatorname{Hom}}^{\otimes,\mathrm{ACU}}(\underline C,\underline Q)$ the category of all $\otimes$-ACU functor from $\underline C$ to $\underline Q$, then there is a canonical equivalence of categories
\[
\resizebox{\linewidth}{!}{$\displaystyle
\underline{\operatorname{Hom}}^{\otimes,\mathrm{ACU}}(\underline C\deleted{\unsure{\otimes|\times}}\unsure{\times|\otimes}\underline C',\underline Q)\longrightarrow
\underline{\operatorname{Hom}}^{\otimes,\mathrm{ACU}}(\underline C,\underline Q)\times
\underline{\operatorname{Hom}}^{\otimes,\mathrm{ACU}}(\underline C',\underline Q)$}
\]
The $\otimes$-ACU category $\underline P$ thus defined is called the \uline{$\otimes$-category of fractions} of the category $\underline C$ with respect to $(\underline C',(F,\breve F))$. The $\otimes$-category of fractions of $\underline C^{\mathrm{is}}$ with respect to $(\underline C^{\mathrm{is}},(\mathrm{id}_{\underline C^{\mathrm{is}}},\mathrm{id}))$ is a Pic-category which is called the \uline{Pic-envelope} of the category $\underline C$, and denoted by $\operatorname{Pic}(\underline C)$.

For an application of the Pic-envelope, we take $\underline C=\unsure{\mathcal P|\mathscr P}(R)$, category of all finitely generated $R$-modules ($R$ a ring) and $\underline P=\operatorname{Pic}(\mathcal P(R))$, then one obtains
\[
\begin{aligned}
\Pi_0(\underline P)&\simeq K^0(R)\\
\Pi_1(\underline P)&\simeq K^1(R)
\end{aligned}
\]
% slanted
\margin{this Pic-envelope is \uline{not} a \uline{strict} Pic-category}
where $K^0(R)$ is the Grothendieck group and $K^1(R)$ the Whitehead group [1]

The use of the $\otimes$-category of fractions of a $\otimes$-ACU category gives us the following result:
% archive sheet number: 56
\scanpage{58}
\mspage{12}
\unsure{let|Let} $\underline C$ be a $\otimes$-ACU category, $Z$ an arbitrary object of $\underline C$ different from the unity object $\underline 1$, $S$ the functor from $\underline C$ to $\underline C$ defined by
\[
X\longmapsto X\otimes Z
\]
The \uline{suspension category} of the $\otimes$-ACU category $\underline C$ defined by the object $Z$ is the triple $(\underline{\unsure{\mathcal P|\mathscr P}},i,\unsure{p|\rho})$ which solves the universal problem for triples $(\underline Q,j,q)$ where $\underline Q$ is a category, $j$ a functor from $\underline C$ to $\underline Q$, and $q$ an equivalence of categories from $\underline Q$ to $\underline Q$, so that the following diagram commutes
\diagram{\begin{tikzcd}[ampersand replacement=\&]
\underline C \arrow[r,"S"] \arrow[d,"j"'] \& \underline C \arrow[d,"j"]\\
\underline Q \arrow[r,"q"] \& \underline Q
\end{tikzcd}}
up to natural isomorphism. In the case where $\underline C$ is the homotopy category of pointed topological spaces $\underline{\mathrm{Htp}}_{\unsure{\times|\ast|x}}$ together with the smash $\wedge$ (the smash $X\wedge Y$ of two spaces $X$ and $Y$, with the base points $x_0$ and $y_0$, is obtained from the product $X\times Y$ by shrinking the subset $X\vee Y=X\times\{y_0\}\cup\{x_0\}\times Y$ to a single point which is taken as the base point of $X\wedge Y$), and the usual ACU constraint; and $Z$ is the 1-sphere $S^1$ hence $S$ is the suspension functor, we get the well-known definition of the suspension category.

Let $\underline C'$ be the $\otimes$-stable subcategory of $\underline C$ \unsure{generatd|generated} by $Z$ and $\underline P$ the $\otimes$-category of fractions of $\underline C$ with respect to $(\underline C',(F,\mathrm{id}))$ where $F:\underline C'\to\underline C$ is the inclusion functor. One obtains a functor $G:\underline{\mathcal P}\to\underline P$ from the suspension category to the $\otimes$-category of fractions $\underline P$. If $G$ is not faithful, that is the case of the homotopy category of pointed topological spaces $\underline{\mathrm{Htp}}_{\ast}$ together with the smash $\wedge$ and the 1-sphere $S^1$; then it is impossible to construct in $\underline{\mathcal P}$ a law $\otimes$ such that $\underline{\mathcal P}$ together with this law is a $\otimes$-ACU category, $iZ$ invertible in $\underline{\mathcal P}$, and $i$ imbedded in a pair $(i,\breve i)$ which is a $\otimes$-ACU functor from $\underline C$ to $\underline{\mathcal P}$.
% archive sheet number: 57
\scanpage{59}
\mspage{13}
\begin{center}References\end{center}

[1] Bass, H. : K-theory and stable algebra. Publ. math. de l'IHES, n$^\circ$ 22,

[2] Bénabou, J. : Thèse, Paris 1966.

[3] Bourbaki : Théorie des ensembles.

[4] \rule[.25ex]{3em}{.45pt} : Algèbre commutative.

[5] \rule[.25ex]{3em}{.45pt} : Algèbre multilinéaire.

[6] Deligne, P. : Champs de Picard strictement commutatifs. SGA 4 $\overline{\underline{\mathrm{XVIII}}}$

[7] Eilenberg, S. and Kelly, G.M : Closed categor\deleted{\unsure{y}}ies. Proceedings of the conference on categorical algebra (421-561). Springer-Verlag 1965.

[8] Freyd, P. : Stable homotopy. \unsure{Procedings|Proceedings} of the conference on categorical algebra (121-176). Springer-Verlag 1965.

[9] Grothendieck, A. : Biextensions de faisceaux de groupes. SGA 7 $\overline{\underline{\mathrm{VII}}}$.

[10] \rule[.25ex]{2.2em}{.45pt} : Catégories cofibrées additives et complexe cotangent relatif. Lecture notes in mathematics N$^\circ$ 79, Springer-Verlag 1968.

[11] \unsure{Mae-Lane|Mac-Lane}, S. : Homology. Springer-Verlag 1967.

[12] \rule[.25ex]{2.3em}{.45pt} : Categorical algebra. Bull\unsure{,|.} Amer. Mat. Soc. 71 (1965)(40-106)

[13] Mitchell, B. : Theory of categories, Academic Press 1965.

[14] Neantro Saavedra Rivano : Thèse, Paris (1970 ?)

[15] \rule[.25ex]{2em}{.45pt} : Catégories tanakiennes. Lecture notes in mathematics N$^\circ$ 265. Springer-Verlag 1972.

[16] Spanier, E : Algebraic topology. Mc Graw-Hill Inc. 1966,
% archive sheet number: 58
